\ArithmeticHead{385}{R}{SUMMATION EINER DIVERGIRENDEN REIHE.}
\begingroup
\small
\makeatletter\tagsleft@true\let\veqno\@@leqno\makeatother
\begin{center}BEMERKUNGEN.\end{center}
\noindent\hspace*{\parindent}Mit der Aufgabe, die »Summe« der divergenten Reihe
\begin{equation}
1-1+2!-3!+4!-5!+{.}\,{.}\,{.}\tag*{(\(\alpha\))}
\end{equation}
\noindent zu finden, hat sich \textsc{Euler} in seiner Abhandlung \textit{De seriebus divergentibus}\textsuperscript{**}) beschäftigt. Die Reihe ist, wie \textsc{Euler} (a. a. O. S. 213, § 13) bemerkt, »a \textsc{Wallisio} hypergeometrica dicta, signis alternantibus instructa«\textsuperscript{***}). Im art. 1. des vorstehenden Bruchstücks schließt sich \textsc{Gauss} ganz an \textsc{Euler} an. \textsc{Euler} bildet (a. a. O. § 19, S. 220) die Reihe [1], zeigt, daß sie formal der Differentialgleichung [2] genügt, stellt dann die allgemeine Lösung von [1] in der Form [11] dar und erklärt die »Summe« von \((\alpha)\) durch die Gleichung
\begin{equation}
c=1-1+2-6+{.}\,{.}\,{.}=e\int_0^1e^{-\frac1x}\frac{dx}{x}\mathpunct{\raisebox{.45ex}{.}}\tag*{(\(\beta\))}
\end{equation}
\noindent Durch Verwandlung der Reihe [1] in einen Kettenbruch (a. a. O. § 21, S. 224) findet \textsc{Euler} dann für \(c\) den Wert \(0{,}5963473621237\)\textsuperscript{\(\dagger\)}). Im Sinne von H. \textsc{Poincaré} kann man die Beziehung der divergenten
\par\medskip\noindent\rule{.24\linewidth}{.4pt}\par\smallskip
{\footnotesize
\noindent\hspace*{\parindent}\textsuperscript{**}) Novi Commentarii Acad. Petropol. \textbf{5} (1754/5) 1760, S. 205
\par\noindent\hspace*{\parindent}\textsuperscript{***}) Hypergeometrisch heißt nach J. \textsc{Wallis} (\textit{Arithmetica infinitorum}, 1655, Propositio 190) eine Reihe, deren \(n\)-tes Glied das Produkt der \(n\) ersten Glieder einer arithmetischen Reihe ist. Auch \textsc{Gauss} bedient sich im art. 29 der Abhandlung \textit{circa seriem}, Werke III, S. 159 dieses Kunstausdrucks.
\par\noindent\hspace*{\parindent}\textsuperscript{\(\dagger\)}) In den \textit{Institutiones calculi differentialis}, 1755, Pars posterior, Cap. I, § 10, III (L. \textsc{Euleri} Opera
\par}
\par\smallskip\noindent X1.\hfill49
\clearpage\SourceLeaf{393}\ArithmeticHead{386}{L}{ANALYSIS.\quad NACHLASS.}
\noindent Potenzreihe [1] zu der Differentialgleichung [2], der sie formal genügt, so aussprechen, daß der Ausdruck
\[
\frac1{x^n}\bigl(\xi_0-x+1!x^2-2!x^3+{.}\,{.}\,{.}+(-1)^n(n-1)!x^n\bigr),
\]
\noindent wo
\begin{equation}
\xi_0=e^{\frac1x}\int_0^x e^{-\frac1x}\frac{dx}{x}\tag*{(\(\gamma\))}
\end{equation}
\noindent ist, für jeden positiven ganzzahligen Wert von \(n\) sich dem Grenzwerte Null nähert, wenn \(x\) als reale positive Größe in den Punkt \(x=0\) einrückt\textsuperscript{*}); man sagt die Reihe [1] stelle die Lösung \((\gamma)\) der Differentialgleichung [2] für kleine positive Werte von \(x\) asymptotisch dar.

\indent\textsc{Gauss} schlägt zur Bestimmung der durch die Gleichung \((\beta)\) definierten Konstanten \(c\) in den artt. 2.—5. der vorstehenden Aufzeichnung den folgenden Weg ein. Er entwickelt die allgemeine Lösung der Differentialgleichung [2] nach Potenzen von \(\omega=1-x\); diese in der Umgebung der regulären Stelle \(x=1\), d. h. also für \(|\omega|<1\) konvergente Entwicklung [3] hängt linear von dem Werte \(c\) ab, den die Lösung für \(x=1\) annimmt, und zwar ist in [3] die mit \(c\) multiplizierte Reihe die Entwicklung von \(e^{\frac1x-1}\), die in der zweiten Zeile in der Klammer stehende Reihe die Entwicklung von
\[
-e^{\frac1x}\int_1^x e^{-\frac1x}\frac{dx}{x}
\]
\noindent nach Potenzen von \(\omega\). In [4] stellt also der Quotient rechts vom Gleichheitszeichen die Funktion
\begin{equation}
-e\int_1^x e^{-\frac1x}\frac{dx}{x}=e\int_0^\omega e^{-\frac1{1-\omega}}\frac{d\omega}{1-\omega}\tag*{(\(\delta\))}
\end{equation}
\par\medskip\noindent\rule{.24\linewidth}{.4pt}\par\smallskip
{\footnotesize
\noindent omnia, ser. I, vol. 10, S. 225) betrachtet \textsc{Euler} die divergente Reihe
\begin{equation}
1-2!+3!-4!+{.}\,{.}\,{.}\tag*{(\(\alpha'\))}
\end{equation}
\noindent und gibt ihre »Summe« mit \(0{,}4036524077\) an, was in den vier letzten Stellen von dem Werte \(1-c\) abweicht. Daß die letzten Stellen des von \textsc{Euler} im \textit{Calculus differentialis} angegebenen Wertes der »Summe« von \((\alpha')\) unrichtig sind, bemerkt \textsc{Mascheroni} in seinen oben S. 369, Fußnote, angeführten \textit{Adnotationes} siehe a. a. O. S. 426), wo er zwischen der »Summe« \(c\) der Reihe \((\alpha)\) und der \textsc{Euler-Mascheroni}schen Konstanten \(A=0{,}57721\,{.}\,{.}\,{.}\) die Beziehung
\[
c+e\left(A-1+\frac1{2{.}2!}-\frac1{3{.}3!}+\frac1{4{.}4!}-{.}\,{.}\right)=0
\]
\noindent herleitet. Die Lösung [11] der Differentialgleichung [2] verwandelt sich durch die Transformation \(x=-\dfrac1{\log z}\), wenn man noch \(0\) als untere Grenze des Integrals nimmt, in \(-\dfrac1z\operatorname{li}z\), die Reihe [1] geht durch dieselbe Transformation in die von \textsc{Mascheroni} a. a. O. S. 426 in der Gleichung (3) angegebene divergente Entwicklung des Integrallogarithmus über.
\par\noindent\hspace*{\parindent}\textsuperscript{*}) Vergl. J. \textsc{Horn}, \textit{Über das Verhalten der Integrale einer linearen Differentialgleichung erster Ordnung in der Umgebung einer Unbestimmtheitsstelle}, \textsc{Crelles} Journal für Mathematik \textbf{120}, 1899, S. 1, wo S. 17 gerade die Differentialgleichung [2] als Beispiel behandelt wird.
\par}
\clearpage\SourceLeaf{394}\ArithmeticHead{387}{R}{BEMERKUNGEN ZUR REIHENLEHRE.\quad ABSCHNITT I.}
\noindent dar und in [5] steht rechts vom Gleichheitszeichen die Entwicklung dieser Funktion in der Umgebung von \(x=1\) oder \(\omega=0\), nach Unterdrückung des Faktors \(\omega\).

\indent Die Funktion \((\delta)\) nähert sich dem Grenzwerte \(c\), wenn \(x\) mit einem zwischen \(-\dfrac\pi2\) und \(\dfrac\pi2\) gelegenen Argument in den Punkt \(0\) einrückt, also jedenfalls wenn die Annäherung zwischen zwei von \(x=0\) auslaufenden Sehnen des Konvergenzkreises der in [5] auftretenden Reihe erfolgt. Die Angabe von \textsc{Gauss}, daß diese Reihe für \(\omega=1\) oder \(x=0\) die durch die Gleichung \((\beta)\) definierte Konstante \(c\) darstellt, wird also\textsuperscript{*}) bestätigt sein, wenn es gelingt, die Konvergenz jener Reihe für \(\omega=1\) zu erweisen. \textsc{Gauss} hat in der hier vorliegenden Aufzeichnung diesen Konvergenzbeweis nicht geliefert, er läßt sich aber in folgender Weise erbringen.

\indent Wie L. \textsc{Fejér} gezeigt hat\textsuperscript{**}) gilt für die Koeffizienten der Entwicklung
\[
\frac{e{.}\,e^{-\frac1{1-\omega}}}{1-\omega}=\gamma_0+\gamma_1\omega+\gamma_2\omega^2+{.}\,{.}
\]
\noindent die asymptotische Darstellung
\[
\gamma_n=\sqrt{\frac e\pi}\,\frac{\sin\left(2\surd n+\frac\pi4\right)+\frac{\lambda(n)}{n^{\frac14}}}{n^{\frac14}},\qquad(n=1,2,3,{.}\,{.}\,{.}),
\]
\noindent wo \(\lambda(n)\) dem absoluten Betrage nach für jedes \(n\) kleiner ist als eine angebbare Konstante. Hieraus ergeben sich für die Entwicklung der Funktion \((\delta)\)
\[
e\int_0^\omega e^{-\frac1{1-\omega}}\frac{d\omega}{1-\omega}=g_1\omega+g_2\omega^2+g_3\omega^3+{.}\,{.}\,{.}
\]
\noindent zwei Folgerungen, die wir einer brieflichen Mitteilung von L. \textsc{Fejér} entnehmen. »Erstens ist
\[
|g_n|=\left|\frac{\gamma_{n-1}}n\right|<\frac{\text{const.}}{n^{\frac54}},
\]
\noindent was die unbedingte Konvergenz der Reihe
\[
g_1+g_2+g_3+{.}\,{.}\,{.},
\]
\noindent d. h. also der Reihe [5] für \(\omega=1\) erkennen läßt. Zweitens zeigt die asymptotische Darstellung von \(\gamma_n\), daß die \(g_n\) in unendlich viele Gruppen zerfallen, in denen abwechselnd das positive bezw. negative Vorzeichen herrscht«. —

\indent Der letztere Umstand erklärt es, weshalb \textsc{Gauss} bei der von ihm (am Ende des art. 5, S. 383) vorgenommenen Summation erster Glieder eine schlechte Annäherung an den wahren Wert der Konstanten \(c\) erhielt.
\par\medskip\noindent\rule{.24\linewidth}{.4pt}\par\smallskip
{\footnotesize
\noindent\hspace*{\parindent}\textsuperscript{*}) Nach einem bekannten Satze von \textsc{Abel}, siehe Oeuvres, nouv. édit. I, Christiania 1881, S. 223, Théorème IV, vergl. A. \textsc{Pringsheim}, Münchener Sitzungsberichte \textbf{27}, 1897, S. 347.
\par\noindent\hspace*{\parindent}\textsuperscript{**}) Siehe Comptes Rendus \textbf{147}, Paris 1908, S. 1040; der Beweis, der sich auf ein von \textsc{Riemann} in seiner Habilitationsschrift über trigonometrische Reihen (§ 13, \textsc{Riemanns} Werke, 2. Aufl., S. 260 ff.) gegebenes Verfahren gründet, ist ausgeführt im Mathem. und naturw. Anzeiger (Értesitö) der ungar. Akademie der Wissensch. \textbf{27}, 1909, S. 1, einen andern Beweis der \textsc{Fejér}schen Formel findet man bei O. \textsc{Perron}, Archiv der Mathem. und Physik, 3. Reihe, \textbf{22}, 1914, S. 329.
\par}
\begin{flushright}49\textsuperscript{*}\end{flushright}
\clearpage\SourceLeaf{395}\ArithmeticHead{388}{L}{ANALYSIS.\quad NACHLASS.}
\indent Im art. [6.] hat nun \textsc{Gauss} versucht, die Beziehung der divergenten Reihe [1] zu der Differentialgleichung [2] in der Weise klar zu stellen, daß er einen Zusammenhang mit den Untersuchungen von \textsc{Laplace} über die homogene lineare Differentialgleichung mit linearen Koeffizienten (siehe die Gleichung [6]) herstellte\textsuperscript{*}). Die Formel [7] weist deutlich auf die asymptotischen Reihenansätze von \textsc{Laplace} hin. \textsc{Gauss} geht durch Differentiation von der Differentialgleichung [2] oder [8] zu der homogenen linearen Differentialgleichung zweiter Ordnung [9] über, die, wenn man \(\dfrac1x\) als neue unabhängige Veränderliche einführt, die Form
\begin{equation}
x\frac{d^2y}{dx^2}+(1-x)\frac{dy}{dx}-y=0\tag*{\([9]'\)}
\end{equation}
\noindent annimmt. Diese Gleichung \([9]'\) ist in der Tat ein besonderer Fall der \textsc{Laplace}schen Differentialgleichung [6]; sie ist zwar in dem \textsc{Gauss}schen Texte nicht ausdrücklich angegeben, aber der Ansatz [10], der ebenso wie [7] zu der asymptotischen Darstellung in der Nähe von \(x=\infty\) gehört, bezieht sich nicht auf [9] sondern auf \([9]'\).

\indent Dem die artt. 1.—[6.] umfassenden Bruchstück haben wir in den artt. [7.] und [8.] zwei andere Aufzeichnungen hinzugefügt, die namentlich mit den Betrachtungen des art. [6.] eng zusammenhängen. Die in [7.] eingeführten Reihen \(\varphi\) und \(\psi\) sind sogenannte \textsc{Bessel}sche oder Zylinderfunktionen, und zwar haben wir in der jetzt üblichen Bezeichnungsweise
\[
\varphi=J_0(2\sqrt{x}),\quad\psi=-x\varphi'=-\sqrt{x}\,J'_0(2\sqrt{x})=\sqrt{x}\,J_1(2\sqrt{x}),
\]
\noindent wo
\[
\begin{aligned}
J_0(t)&=\sum_{n=0}^\infty\frac{(-1)^n}{(n!)^2}\left(\frac t2\right)^{2n},\\[.5em]
J_1(t)&=\sum_{n=0}^\infty\frac{(-1)^n}{n!(n+1)!}\left(\frac t2\right)^{2n+1}
\end{aligned}
\]
\noindent ist\textsuperscript{**}). Die Differentialgleichung [12], der \(\varphi\) genügt, ist ein besonderer Fall der \textsc{Laplace}schen Gleichung [6] des art. [6.]. Daß \textsc{Gauss} die für große Werte von \(x\) geltenden asymptotischen Darstellungen der ganzen transzendenten Funktionen \(\varphi\) und \(\psi\) betrachtet hat, zeigt die im art. [8.] wiedergegebene Aufzeichnung Nr. 82 des \textit{Tagebuchs}, wo in der Gleichung [17] eine asymptotische Bestimmung für die Nullstellen der Funktion \(\psi\) gegeben wird\textsuperscript{***}) J. \textsc{Horn} bestätigt die Gleichung [17] durch die folgende Rechnung.
\par\medskip\noindent\rule{.24\linewidth}{.4pt}\par\smallskip
{\footnotesize
\noindent\hspace*{\parindent}\textsuperscript{*}) Siehe \textsc{de la Place}, \textit{Mémoire sur les approximations des formules qui sont fonctions de très-grands nombres}, Histoire de l’Académie des Sciences, année 1782, Paris 1785, Mémoires, S. 1 ff.; siehe insbesondere S. 47 ff.
\par\noindent\hspace*{\parindent}\textsuperscript{**}) Die Funktion \(J_0(t)\) findet sich wohl zuerst bei \textsc{Daniel Bernoulli}, \textit{Theoremata de oscillationibus corporum etc.} Commentarii Acad. Petrop. \textbf{6} (1732/33) 1738, S. 116, 118, die Differentialgleichung, der sie Genüge leistet, ebenda \textbf{7} (1734/35) 1740, S. 171. Die allgemeine \textsc{Bessel}sche Funktion \(J_\nu(t)\) kommt in der Abhandlung von \textsc{Euler} \textit{De motu vibratorio tympanorum}, Novi Commentarii Acad. Petropol. \textbf{10}, (1764) 1766, S. 243, auf S. 256 vor; diese Abhandlung hat \textsc{Gauss} jedenfalls gekannt.
\par\noindent\hspace*{\parindent}\textsuperscript{***}) Die asymptotische Darstellung von \(J_0(t)\) findet sich zuerst veröffentlicht bei S. D. \textsc{Poisson}, \textit{Sur la distribution de la chaleur dans les corps solides}, Journal de l’École Polytechnique, cah. 19, 1823, S. 249, auf S. 350; die analoge Darstellung für \(J_1(t)\) gibt wohl zum ersten Male P. A. \textsc{Hansen}, \textit{Ermittelung der absoluten Störungen u.s.w.} I. Theil, Schriften der Sternwarte Seeberg, Gotha, 1843, S. 115, § 51, der auch in einer Fußnote auf S. 116 die asymptotische Formel
\par}
\clearpage\SourceLeaf{396}\ArithmeticHead{389}{R}{BEMERKUNGEN ZUR REIHENLEHRE.\quad ABSCHNITT I.}
\indent »Für die \textsc{Bessel}sche Funktion \(J_1(x)\) gilt\textsuperscript{*}) die asymptotische Darstellung
\[
J_1(x)=\sqrt{\frac2{\pi x}}\left\{P\cos\left(\frac{3\pi}4-x\right)+Q\sin\left(\frac{3\pi}4-x\right)\right\},
\]
\noindent wo
\[
P=1+\frac{15}{128x^2}-{.}\,{.}\,{.},\quad Q=\frac3{8x}-\frac{105}{1024x^3}+{.}\,{.}\,{.}
\]
\noindent ist. Die Gleichung \(J_1(x)=0\) schreibt sich also
\[
\operatorname{cotg}\left(\frac{3\pi}4-x\right)=\operatorname{tg}\left(x-\left(k+\frac14\right)\pi\right)=-\frac QP=-\frac3{8x}+\frac{75}{512x^3}\,{.}\,{.}\,{.},
\]
\noindent wo \(k\) eine ganze Zahl bedeutet; hieraus folgt
\[
\left(k+\frac14\right)\pi-x=\operatorname{arctg}\left(\frac3{8x}-\frac{75}{512x^3}\,{.}\,{.}\,{.}\right)=\frac3{8x}-\frac{21}{128x^3}\,{.}\,{.}\,{.},
\]
\noindent also
\[
x+\frac3{8x}-\frac{21}{128x^3}\,{.}\,{.}\,{.}=\left(k+\frac14\right)\pi.
\]
\noindent Aus der Gleichung \(\psi=0\) oder \(J_1(2\sqrt{x})=0\) folgt demnach
\[
2\sqrt{x}+\frac3{16\sqrt{x}}-\frac{21}{1024x\sqrt{x}}\,{.}\,{.}\,{.}=\left(k+\frac14\right)\pi,
\]
\noindent was mit der Formel [17] übereinstimmt«.

\indent Es steht hiernach fest, daß \textsc{Gauss} die asymptotische Darstellung der Funktion \(\psi\), mindestens so weit, als sie in der vorstehenden Rechnung benutzt wird, am 16 Oktober 1797, wo er die Tagebuchnotiz Nr. 82 schrieb, gekannt hat. Aus derselben Zeit dürften auch die Leisteaufzeichnung art. [7.] und der die artt. 1.—[6.] enthaltende Zettel stammen. Untersuchungen zum \textsc{Kepler}schen Problem, die \textsc{Gauss} einige Jahre später angestellt hat, und die mit den hier besprochenen Ansätzen zusammenhängen, sind weiter unten im Abschnitt [VII.] abgedruckt; vergl. die Bemerkungen zum Abschnitt [VII.] und zu dem auf diesen Abschnitt folgenden Briefwechsel.
\begin{flushright}\textsc{Schlesinger}.\end{flushright}
\par\medskip\noindent\rule{.24\linewidth}{.4pt}\par\smallskip
{\footnotesize
\[
\lambda=\frac12\pi\left(k+\frac12\right)-\frac3{32\lambda}+\frac{21}{2048\lambda^3}-\frac{1899}{327680\lambda^5}\pm\text{etc.}
\]
\noindent für die Wurzeln der Gleichung \(J_1(t)=0\) aufstellt. Wie man sieht, stimmt diese Formel mit der \textsc{Gauss}schen Gleichung [17] überein.
\par\noindent\hspace*{\parindent}\textsuperscript{*}) Siehe \textsc{Hansen} a. a. O.
\par}
\ClosingRule
\clearpage\endgroup
